\section{Methodology}
\label{sec:methodology}

Building on our observation that duality provides a principled mapping from attention to SSM parameters, we design a two-stage verification protocol that isolates parameter validity from structural fidelity.

\subsection{Overview}

Our methodology tests whether Mamba-2 duality equations can initialize SSM parameters from BERT attention weights such that:

\begin{enumerate}
    \item \textbf{Stage 1 (Existence - h-e1):} The derived parameters produce stable, non-divergent SSM outputs.
    \item \textbf{Stage 2 (Mechanism - h-m1):} The derived parameters reconstruct attention outputs better than random initialization.
\end{enumerate}

\subsection{Duality-Based Parameter Extraction}

We implement closed-form SSM parameter derivation from BERT attention weights following the Mamba-2 SSD framework. For each attention layer with query, key, and value projections ($W_Q, W_K, W_V \in \mathbb{R}^{768 \times 768}$):

\textbf{A Matrix (State Dynamics):} We compute $QK^T = W_Q \cdot W_K^T$ and extract its principal components via SVD. The negative absolute values ensure stability (negative eigenvalues prevent divergence).

\textbf{B and C Matrices:} We derive $B$ and $C$ from the SVD components, normalized for numerical stability.

\textbf{D and $\Delta$:} Skip connection $D = 0.1$ and discretization step $\Delta = 1/\sqrt{768}$.

\subsection{Stability Validation (h-e1)}

We validate parameter stability by running SSM forward passes on calibration data, checking for NaN/Inf and measuring magnitude ratio relative to Transformer outputs.

\subsection{Reconstruction Error Comparison (h-m1)}

We compare duality initialization against random initialization using output Frobenius norm: $\|SSM\ output - Attention\ output\|_F$.
